% Auto-generated complete chapter from 08B_twinsight_case_study_final_structure.md
\chapter{08B\_twinsight\_case\_study\_final\_structure}
\label{complete:root_010}

\begin{sourcemeta}
\textbf{Fuente:} \href{file:///E:/Laboral/08B_twinsight_case_study_final_structure.md}{08B\_twinsight\_case\_study\_final\_structure.md}\\
\textbf{Seccion:} root\\
\textbf{Estado:} canonical\\
\textbf{Rol:} Modulo final complementario\\
\textbf{Lineas originales:} 1111\\
\textbf{Ultima modificacion:} 2026-06-11 19:42\\
\textbf{Deduplicacion conservadora:} 0 bloques repetidos omitidos; 0 caracteres repetidos no reimpresos.
\end{sourcemeta}

\section*{Resumen de orientacion}
Use this framing across portfolio, LinkedIn, GitHub and ArtStation:

\section*{Contenido completo depurado}
\addcontentsline{toc}{section}{Contenido completo depurado}




\subsection{0. Module status}




\begin{mdcode}
Bloque de codigo: text
Module: 08B_twinsight_case_study_final_structure.md
Date generated: 2026-06-12
Status: Complete v1
Purpose: Final portfolio case-study structure for TwinSight X500
Inputs: 08, 17, 19, 20, 21, 21B, 29, 29B, 29C, B1, B5
Use: Portfolio page, GitHub README, LinkedIn Featured, ArtStation breakdown, interview follow-up
\end{mdcode}




\medskip\hrule\medskip




\section{1. Final positioning}




\subsection{1.1 Primary project identity}




\begin{mdcode}
Bloque de codigo: text
TwinSight X500
Unity WebGL Technical Visualization for Drone Assembly Inspection
\end{mdcode}




\subsection{1.2 One-line description}




\begin{mdcode}
Bloque de codigo: text
An interactive browser-based 3D viewer that helps users inspect and understand the assembly of a Holybro X500 V2 drone through component selection, exploded view, cross-section tools, visual modes and technical information panels.
\end{mdcode}




\subsection{1.3 Professional framing}




Use this framing across portfolio, LinkedIn, GitHub and ArtStation:




\begin{mdcode}
Bloque de codigo: text
TwinSight X500 is a Unity WebGL technical visualization prototype for drone assembly inspection. It transforms CAD-derived drone assets into an optimized browser-based 3D viewer with component-level interaction, exploded assembly visualization, cross-section inspection, visual modes and formative usability/workload evaluation.
\end{mdcode}




\subsection{1.4 Role-aligned keywords}




\begin{mdcode}
Bloque de codigo: text
Unity
C#
Unity WebGL
Real-Time 3D
Technical Visualization
Technical Art
CAD-to-Realtime
Interactive 3D
URP
UI Toolkit
Shader Graph
Blender
Optimization
Exploded View
Clipping
Cross-Section
SUS
NASA-TLX
Think-Aloud
\end{mdcode}




\subsection{1.5 What this project should prove}




\begin{center}
\scriptsize
\begin{longtable}{>{\raggedright\arraybackslash}p{0.455\linewidth} >{\raggedright\arraybackslash}p{0.455\linewidth}}
\toprule
Proof area & What TwinSight must communicate \\
\midrule
\endfirsthead
\toprule
Proof area & What TwinSight must communicate \\
\midrule
\endhead
Real-time 3D development & interactive browser-based Unity application \\
Technical art & visual modes, optimization, rendering constraints \\
CAD-to-realtime & heavy CAD data converted into optimized runtime asset \\
Technical visualization & assembly understanding, inspection, spatial relationships \\
UI/UX & technical panels, bottom-sheet interaction, mobile-first logic \\
Evaluation & SUS, NASA-TLX and Think-Aloud validation \\
Portfolio maturity & clear case study, video, metrics, limitations, links \\
\bottomrule
\end{longtable}
\end{center}




\medskip\hrule\medskip




\section{2. Page hierarchy}




\subsection{2.1 Recommended case-study page order}




\begin{mdcode}
Bloque de codigo: text
01. Hero
02. Short value proposition
03. Demo video
04. Problem
05. Solution
06. Key features
07. Technical pipeline
08. CAD-to-realtime optimization
09. Unity WebGL architecture
10. Interaction systems
11. Visual modes
12. Evaluation results
13. Role and contribution
14. AI-assisted workflow disclosure
15. Constraints and limitations
16. Screenshots / media gallery
17. Links
18. Next steps
\end{mdcode}




\subsection{2.2 Recruiter-first logic}




The first screen must answer:




\begin{mdcode}
Bloque de codigo: text
What is it?
What role does it support?
What technologies were used?
What can the user do?
Why is it credible?
\end{mdcode}




Do not start with academic methodology.




\medskip\hrule\medskip




\section{3. Hero section}




\subsection{3.1 Hero title}




\begin{mdcode}
Bloque de codigo: text
TwinSight X500
\end{mdcode}




\subsection{3.2 Hero subtitle}




\begin{mdcode}
Bloque de codigo: text
Unity WebGL technical visualization for drone assembly inspection.
\end{mdcode}




\subsection{3.3 Hero paragraph}




\begin{mdcode}
Bloque de codigo: text
TwinSight X500 is an interactive browser-based 3D viewer for inspecting the assembly of a Holybro X500 V2 drone. It was built in Unity WebGL with C#, URP and UI Toolkit, using a CAD-to-realtime pipeline to transform complex source geometry into an optimized real-time visualization.
\end{mdcode}




\subsection{3.4 Hero metadata}




\begin{center}
\scriptsize
\begin{longtable}{>{\raggedright\arraybackslash}p{0.455\linewidth} >{\raggedright\arraybackslash}p{0.455\linewidth}}
\toprule
Field & Value \\
\midrule
\endfirsthead
\toprule
Field & Value \\
\midrule
\endhead
Role & Real-Time 3D Developer / Unity Technical Artist \\
Stack & Unity, C\#, WebGL, URP, UI Toolkit, Blender \\
Domain & Technical visualization / CAD-to-realtime \\
Project type & Academic/portfolio prototype \\
Platform & Browser / WebGL \\
Evaluation & SUS, NASA-TLX Raw, Think-Aloud \\
\bottomrule
\end{longtable}
\end{center}




\subsection{3.5 Hero CTA buttons}




\begin{mdcode}
Bloque de codigo: text
Watch demo
View GitHub
Read technical breakdown
Open WebGL prototype
\end{mdcode}




Only include “Open WebGL prototype” if the build is public and stable.




\medskip\hrule\medskip




\section{4. Demo video placement}




\subsection{4.1 Primary video}




Embed:




\begin{mdcode}
Bloque de codigo: text
twinsight_x500_demo_90s.mp4
\end{mdcode}




or hosted link:




\begin{mdcode}
Bloque de codigo: text
YouTube / Vimeo / portfolio-hosted video
\end{mdcode}




\subsection{4.2 Caption under video}




\begin{mdcode}
Bloque de codigo: text
90-second overview showing browser-based interaction, component selection, exploded view, clipping, visual modes, CAD optimization and evaluation results.
\end{mdcode}




\subsection{4.3 Secondary video}




Optional:




\begin{mdcode}
Bloque de codigo: text
twinsight_x500_technical_breakdown_5min.mp4
\end{mdcode}




Caption:




\begin{mdcode}
Bloque de codigo: text
Technical breakdown covering CAD preparation, Unity WebGL architecture, runtime interaction systems, visual modes, UI Toolkit and evaluation methodology.
\end{mdcode}




\medskip\hrule\medskip




\section{5. Problem section}




\subsection{5.1 Section title}




\begin{mdcode}
Bloque de codigo: text
Problem: 2D documentation makes spatial relationships harder to understand
\end{mdcode}




\subsection{5.2 Final copy}




\begin{mdcode}
Bloque de codigo: text
Drone assembly documentation is often presented through static 2D diagrams, part lists and written instructions. This forces users to mentally reconstruct spatial relationships between components, orientations, fasteners and hidden structures.

For complex assemblies such as a Holybro X500 V2 drone, this increases cognitive effort and makes it harder to understand how parts relate to each other in 3D space.
\end{mdcode}




\subsection{5.3 Problem bullets}




\begin{mdcode}
Bloque de codigo: text
- Static diagrams do not allow direct inspection.
- Hidden structures require mental reconstruction.
- Part relationships are difficult to infer from 2D views.
- Assembly hierarchy is not always visually obvious.
- Users cannot easily isolate, cut through or compare components.
\end{mdcode}




\subsection{5.4 Visuals for problem section}




\begin{center}
\scriptsize
\begin{longtable}{>{\raggedright\arraybackslash}p{0.455\linewidth} >{\raggedright\arraybackslash}p{0.455\linewidth}}
\toprule
Visual & Purpose \\
\midrule
\endfirsthead
\toprule
Visual & Purpose \\
\midrule
\endhead
2D diagram / manual-style screenshot & show limitation \\
dense part list & show complexity \\
drone exploded diagram reference & show assembly complexity \\
side-by-side 2D vs 3D & communicate transition \\
\bottomrule
\end{longtable}
\end{center}




Use only visuals you have rights to use or self-created diagrams.




\medskip\hrule\medskip




\section{6. Solution section}




\subsection{6.1 Section title}




\begin{mdcode}
Bloque de codigo: text
Solution: interactive technical visualization in the browser
\end{mdcode}




\subsection{6.2 Final copy}




\begin{mdcode}
Bloque de codigo: text
TwinSight X500 converts CAD-derived drone geometry into an optimized Unity WebGL viewer. Instead of reading static documentation, users can navigate the assembly, select components, inspect technical information, activate exploded view, cut through the model and switch between visual modes.

The goal is not only to show the drone, but to make the relationships between its parts easier to understand.
\end{mdcode}




\subsection{6.3 Solution bullets}




\begin{mdcode}
Bloque de codigo: text
- Browser-based Unity WebGL viewer.
- Component-level selection and highlighting.
- Technical part information panels.
- Exploded view for assembly hierarchy.
- Cross-section / clipping for hidden structures.
- Multiple visual modes for inspection.
- CAD-to-realtime optimization pipeline.
- Formative usability and workload evaluation.
\end{mdcode}




\subsection{6.4 Solution diagram}




Recommended diagram:




\begin{mdcode}
Bloque de codigo: text
2D documentation
      ↓
CAD-derived 3D source
      ↓
Blender optimization / cleanup
      ↓
Unity WebGL interactive viewer
      ↓
User inspection + evaluation
\end{mdcode}




\medskip\hrule\medskip




\section{7. Key features section}




\subsection{7.1 Feature cards}




\begin{center}
\scriptsize
\begin{longtable}{>{\raggedright\arraybackslash}p{0.305\linewidth} >{\raggedright\arraybackslash}p{0.305\linewidth} >{\raggedright\arraybackslash}p{0.305\linewidth}}
\toprule
Feature & Description & Visual \\
\midrule
\endfirsthead
\toprule
Feature & Description & Visual \\
\midrule
\endhead
Component selection & Select individual parts and inspect related data & selected arm / motor / frame part \\
Exploded view & Separate assemblies to reveal hierarchy and placement & full exploded drone \\
Cross-section tools & Cut through geometry to inspect hidden structures & clipping plane screenshot \\
Visual modes & Switch between realistic, X-ray, ghosted, blueprint, wireframe and thermal-style views & visual modes grid \\
Technical UI & Access part information through contextual panels & UI panel close-up \\
CAD optimization & Convert heavy CAD-derived assets into WebGL-ready geometry & before/after wireframe \\
Evaluation & Validate usability and workload with structured methods & metrics card \\
\bottomrule
\end{longtable}
\end{center}




\subsection{7.2 Feature copy blocks}




\subsubsection{Component selection}




\begin{mdcode}
Bloque de codigo: text
Users can select components directly in the 3D viewer. The selected part is highlighted and linked to contextual technical information, reducing the need to search through separate lists or diagrams.
\end{mdcode}




\subsubsection{Exploded view}




\begin{mdcode}
Bloque de codigo: text
The exploded view separates assemblies while preserving spatial logic. This helps users understand how components relate to each other and where each part belongs in the full structure.
\end{mdcode}




\subsubsection{Cross-section / clipping}




\begin{mdcode}
Bloque de codigo: text
The clipping tool cuts through the model to reveal internal structures and hidden relationships. It supports technical inspection without requiring users to disassemble the visual context.
\end{mdcode}




\subsubsection{Visual modes}




\begin{mdcode}
Bloque de codigo: text
Multiple visual modes support different inspection tasks: realistic view for material context, X-ray and ghosted modes for internal visibility, blueprint and wireframe modes for technical reading, and a thermal-style qualitative mode for visual exploration.
\end{mdcode}




\subsubsection{Technical UI}




\begin{mdcode}
Bloque de codigo: text
The interface provides technical part information while preserving the 3D inspection context. The UI was designed with a mobile-first bottom-sheet logic and runtime interaction in mind.
\end{mdcode}




\medskip\hrule\medskip




\section{8. Technical pipeline section}




\subsection{8.1 Section title}




\begin{mdcode}
Bloque de codigo: text
Technical pipeline: CAD-to-realtime optimization
\end{mdcode}




\subsection{8.2 Final copy}




\begin{mdcode}
Bloque de codigo: text
The source drone model came from CAD-derived geometry that was not suitable for direct WebGL deployment. The pipeline focused on preserving technical readability while reducing geometry complexity, preparing assets for real-time interaction and organizing the model for runtime selection and visualization.
\end{mdcode}




\subsection{8.3 Pipeline stages}




\begin{mdcode}
Bloque de codigo: text
01. CAD source inspection
02. STEP / tessellation route exploration
03. Blender cleanup and reconstruction
04. Modular fastener handling
05. High-poly / low-poly preparation
06. UV and bake preparation
07. Material and visual-mode setup
08. Unity import and scene organization
09. Runtime interaction systems
10. WebGL build and evaluation
\end{mdcode}




\subsection{8.4 Pipeline table}




\begin{center}
\scriptsize
\begin{longtable}{>{\raggedright\arraybackslash}p{0.305\linewidth} >{\raggedright\arraybackslash}p{0.305\linewidth} >{\raggedright\arraybackslash}p{0.305\linewidth}}
\toprule
Stage & Tool & Purpose \\
\midrule
\endfirsthead
\toprule
Stage & Tool & Purpose \\
\midrule
\endhead
CAD source & STEP / CAD route & source geometry \\
Tessellation & MoI3D / Stepper route & generate usable mesh \\
Cleanup & Blender & remove/repair/rebuild geometry \\
Optimization & Blender & reduce geometry for WebGL \\
UV/Bake & RizomUV / Marmoset / Blender & prepare textures if used \\
Runtime & Unity & interaction and visualization \\
UI & UI Toolkit & technical panels and controls \\
Deployment & Unity WebGL & browser delivery \\
Evaluation & SUS / NASA-TLX / Think-Aloud & usability/workload validation \\
\bottomrule
\end{longtable}
\end{center}




Only list tools actually used in final public copy. If RizomUV or Marmoset were planned but not used, mark them as planned or omit them.




\medskip\hrule\medskip




\section{9. Optimization section}




\subsection{9.1 Section title}




\begin{mdcode}
Bloque de codigo: text
Optimization: from heavy CAD data to WebGL-ready geometry
\end{mdcode}




\subsection{9.2 Final copy}




\begin{mdcode}
Bloque de codigo: text
A central technical challenge was converting complex CAD-derived geometry into a model that could run in a browser while preserving the visual information needed for assembly inspection.
\end{mdcode}




\subsection{9.3 Metrics card}




\begin{center}
\scriptsize
\begin{longtable}{>{\raggedright\arraybackslash}p{0.455\linewidth} >{\raggedright\arraybackslash}p{0.455\linewidth}}
\toprule
Metric & Value \\
\midrule
\endfirsthead
\toprule
Metric & Value \\
\midrule
\endhead
Optimized geometry & 95,617 triangles \\
Source CAD routes & 6.5M+ triangles \\
Canonical research parts & 28 \\
Scene nodes / anchors & 30 \\
Technical elements & 257 \\
Evaluation participants & 12 \\
Task-condition records & 96 \\
\bottomrule
\end{longtable}
\end{center}




\subsection{9.4 Safe metric language}




Use:




\begin{mdcode}
Bloque de codigo: text
Reduced CAD-derived source routes exceeding 6.5M triangles into a 95,617-triangle optimized model for browser-based real-time inspection.
\end{mdcode}




Avoid:




\begin{mdcode}
Bloque de codigo: text
Optimized the model by 98.5% and proved production scalability.
\end{mdcode}




The first is factual. The second overclaims.




\subsection{9.5 Required visuals}




\begin{center}
\scriptsize
\begin{longtable}{>{\raggedright\arraybackslash}p{0.455\linewidth} >{\raggedright\arraybackslash}p{0.455\linewidth}}
\toprule
Visual & Required \\
\midrule
\endfirsthead
\toprule
Visual & Required \\
\midrule
\endhead
high-poly / source route screenshot & Yes \\
optimized model screenshot & Yes \\
wireframe comparison & Yes \\
triangle metric card & Yes \\
WebGL browser proof & Yes \\
Blender cleanup screenshot & Recommended \\
Unity profiler screenshot & Recommended \\
\bottomrule
\end{longtable}
\end{center}




\medskip\hrule\medskip




\section{10. Unity WebGL architecture section}




\subsection{10.1 Section title}




\begin{mdcode}
Bloque de codigo: text
Unity WebGL architecture
\end{mdcode}




\subsection{10.2 Final copy}




\begin{mdcode}
Bloque de codigo: text
The Unity application coordinates runtime selection, visualization states, UI panels and inspection tools inside a WebGL build. The architecture separates model organization, user interaction, visual states and interface logic so that the viewer can support multiple inspection modes without becoming a static 3D model viewer.
\end{mdcode}




\subsection{10.3 Architecture components}




\begin{mdcode}
Bloque de codigo: text
- Model hierarchy and selectable parts
- Selection and highlight system
- Camera and navigation logic
- Exploded-view state controller
- Clipping / cross-section controller
- Visual-mode switching system
- Technical metadata / part information
- UI Toolkit panels
- WebGL deployment configuration
\end{mdcode}




\subsection{10.4 Architecture visual}




Recommended diagram:




\begin{mdcode}
Bloque de codigo: text
User input
   ↓
Interaction controller
   ↓
Selection / camera / clipping / exploded states
   ↓
Visualization systems
   ↓
UI Toolkit panels + 3D model feedback
   ↓
WebGL runtime
\end{mdcode}




\medskip\hrule\medskip




\section{11. Visual modes section}




\subsection{11.1 Section title}




\begin{mdcode}
Bloque de codigo: text
Visual modes for technical inspection
\end{mdcode}




\subsection{11.2 Visual mode table}




\begin{center}
\scriptsize
\begin{longtable}{>{\raggedright\arraybackslash}p{0.305\linewidth} >{\raggedright\arraybackslash}p{0.305\linewidth} >{\raggedright\arraybackslash}p{0.305\linewidth}}
\toprule
Mode & Purpose & Claim safety \\
\midrule
\endfirsthead
\toprule
Mode & Purpose & Claim safety \\
\midrule
\endhead
Realistic & material and form readability & safe \\
X-ray & internal visibility & safe if implemented \\
Ghosted & semi-transparent inspection & safe \\
Blueprint & technical drawing aesthetic & safe \\
Wireframe & structure/topology reading & safe \\
Solid & neutral inspection & safe \\
Thermal-style & qualitative visual exploration & must clarify not simulation \\
\bottomrule
\end{longtable}
\end{center}




\subsection{11.3 Required disclaimer}




\begin{mdcode}
Bloque de codigo: text
The thermal-style mode is a qualitative visual mode. It is not a physical thermal simulation, FEA result or sensor-driven thermal analysis.
\end{mdcode}




\subsection{11.4 Visual layout}




Use a grid:




\begin{mdcode}
Bloque de codigo: text
Realistic | X-ray | Ghosted
Blueprint | Wireframe | Solid
Thermal-style | Clipping | Exploded
\end{mdcode}




\medskip\hrule\medskip




\section{12. Evaluation section}




\subsection{12.1 Section title}




\begin{mdcode}
Bloque de codigo: text
Evaluation: usability and perceived workload
\end{mdcode}




\subsection{12.2 Final copy}




\begin{mdcode}
Bloque de codigo: text
TwinSight X500 was evaluated as an academic/portfolio prototype using usability, workload and qualitative observation methods. The goal was to compare how users interacted with the 3D viewer and traditional 2D support during assembly-understanding tasks.
\end{mdcode}




\subsection{12.3 Evaluation methods}




\begin{center}
\scriptsize
\begin{longtable}{>{\raggedright\arraybackslash}p{0.455\linewidth} >{\raggedright\arraybackslash}p{0.455\linewidth}}
\toprule
Method & Purpose \\
\midrule
\endfirsthead
\toprule
Method & Purpose \\
\midrule
\endhead
SUS & usability perception \\
NASA-TLX Raw & perceived workload \\
Think-Aloud & qualitative insight \\
Task-condition records & structured interaction data \\
\bottomrule
\end{longtable}
\end{center}




\subsection{12.4 Evaluation metrics}




\begin{center}
\scriptsize
\begin{longtable}{>{\raggedright\arraybackslash}p{0.455\linewidth} >{\raggedright\arraybackslash}p{0.455\linewidth}}
\toprule
Metric & Value \\
\midrule
\endfirsthead
\toprule
Metric & Value \\
\midrule
\endhead
Participants & 12 \\
SUS average & 91.88 \\
NASA-TLX Raw, 3D viewer & 8.69 \\
NASA-TLX Raw, 2D support & 19.89 \\
Task-condition records & 96 \\
\bottomrule
\end{longtable}
\end{center}




\subsection{12.5 Safe interpretation}




Use:




\begin{mdcode}
Bloque de codigo: text
The formative evaluation suggested high perceived usability and lower perceived workload for the 3D viewer condition compared with 2D support in the tested tasks.
\end{mdcode}




Avoid:




\begin{mdcode}
Bloque de codigo: text
TwinSight proves that 3D documentation is always better than 2D documentation.
\end{mdcode}




\subsection{12.6 Evaluation visuals}




\begin{center}
\scriptsize
\begin{longtable}{>{\raggedright\arraybackslash}p{0.455\linewidth} >{\raggedright\arraybackslash}p{0.455\linewidth}}
\toprule
Visual & Use \\
\midrule
\endfirsthead
\toprule
Visual & Use \\
\midrule
\endhead
SUS metric card & recruiter-readable proof \\
NASA-TLX comparison & workload evidence \\
participant/task count & credibility \\
Think-Aloud quote snippets & qualitative support \\
evaluation flow diagram & academic grounding \\
\bottomrule
\end{longtable}
\end{center}




\medskip\hrule\medskip




\section{13. Role and contribution section}




\subsection{13.1 Section title}




\begin{mdcode}
Bloque de codigo: text
My role
\end{mdcode}




\subsection{13.2 Final copy}




\begin{mdcode}
Bloque de codigo: text
I worked across the complete prototype pipeline: CAD-to-realtime asset preparation, Unity WebGL implementation, C# runtime interaction systems, technical UI, visual modes, documentation and formative evaluation.
\end{mdcode}




\subsection{13.3 Contribution bullets}




\begin{mdcode}
Bloque de codigo: text
- Prepared and optimized CAD-derived drone geometry for real-time WebGL use.
- Organized selectable components and scene hierarchy.
- Implemented runtime interaction systems in Unity/C#.
- Built component selection, highlighting and information panels.
- Developed exploded view and clipping/cross-section interaction.
- Implemented multiple visual inspection modes.
- Designed the technical UI and mobile-first interaction logic.
- Prepared evaluation workflow using SUS, NASA-TLX Raw and Think-Aloud.
- Documented constraints, findings and future improvements.
\end{mdcode}




\subsection{13.4 Role clarity note}




If the thesis project involved AI assistance, external assets or source CAD files, clarify the boundary.




Use:




\begin{mdcode}
Bloque de codigo: text
The source CAD model was used as the technical basis for the drone geometry. My work focused on transforming, optimizing, organizing and implementing it as an interactive real-time WebGL visualization prototype.
\end{mdcode}




\medskip\hrule\medskip




\section{14. AI-assisted workflow disclosure}




\subsection{14.1 Why include this}




AI assistance can become a trust risk if hidden and discovered later.




The disclosure should be clear, controlled and not self-sabotaging.




\subsection{14.2 Recommended disclosure}




\begin{mdcode}
Bloque de codigo: text
AI assistance was used as an implementation and iteration support tool during parts of the development process. Final architecture decisions, debugging, Unity integration, asset preparation, evaluation design, documentation and project ownership remained my responsibility.
\end{mdcode}




\subsection{14.3 Short version}




\begin{mdcode}
Bloque de codigo: text
AI-assisted implementation support; final integration, validation and technical ownership by Alexander Woodcock Salomón.
\end{mdcode}




\subsection{14.4 What not to say}




Avoid:




\begin{mdcode}
Bloque de codigo: text
AI built the project.
I used ChatGPT for everything.
I do not know how parts of it work.
\end{mdcode}




If asked in interview, answer precisely and separate:




\begin{mdcode}
Bloque de codigo: text
assisted coding
debugging
architecture
integration
validation
final ownership
\end{mdcode}




\medskip\hrule\medskip




\section{15. Limitations section}




\subsection{15.1 Section title}




\begin{mdcode}
Bloque de codigo: text
Constraints and limitations
\end{mdcode}




\subsection{15.2 Final copy}




\begin{mdcode}
Bloque de codigo: text
TwinSight X500 is an academic/portfolio prototype, not a deployed industrial product. It does not include live sensor integration, predictive maintenance, WebAR, multi-user collaboration or production monitoring. The current focus is technical visualization and assembly understanding through browser-based 3D interaction.
\end{mdcode}




\subsection{15.3 Limitation table}




\begin{center}
\scriptsize
\begin{longtable}{>{\raggedright\arraybackslash}p{0.455\linewidth} >{\raggedright\arraybackslash}p{0.455\linewidth}}
\toprule
Limitation & Status \\
\midrule
\endfirsthead
\toprule
Limitation & Status \\
\midrule
\endhead
Not a production digital twin & true \\
No live IoT data & true \\
No predictive maintenance & true \\
No WebAR placement & true \\
No real thermal simulation & true \\
Academic participant sample & true \\
WebGL constraints & true \\
Prototype scope & true \\
\bottomrule
\end{longtable}
\end{center}




\subsection{15.4 Why this helps}




Limitations increase credibility. They prevent overclaiming and make the project easier to defend in interviews.




\medskip\hrule\medskip




\section{16. Next steps section}




\subsection{16.1 Section title}




\begin{mdcode}
Bloque de codigo: text
Future improvements
\end{mdcode}




\subsection{16.2 Safe future improvements}




\begin{mdcode}
Bloque de codigo: text
- Add guided assembly sequences.
- Improve metadata depth for each component.
- Add search and part filtering.
- Add WebGPU or performance-focused rendering experiments.
- Add optional XR inspection mode.
- Add sensor/state data as a digital-twin extension.
- Expand to other drones or technical assemblies.
- Improve mobile interaction and loading performance.
\end{mdcode}




\subsection{16.3 Do not present future work as current functionality}




Use:




\begin{mdcode}
Bloque de codigo: text
future extension
planned improvement
possible next step
\end{mdcode}




Do not use:




\begin{mdcode}
Bloque de codigo: text
includes
supports
delivers
\end{mdcode}




unless implemented.




\medskip\hrule\medskip




\section{17. Required screenshots and media}




\subsection{17.1 Minimum media package}




\begin{mdcode}
Bloque de codigo: text
01_hero_final_view.jpg
02_browser_webgl_view.jpg
03_component_selection.jpg
04_technical_panel.jpg
05_exploded_view.jpg
06_cross_section.jpg
07_visual_modes_grid.jpg
08_cad_optimization_before_after.jpg
09_wireframe_comparison.jpg
10_metrics_card.jpg
11_pipeline_diagram.jpg
12_architecture_diagram.jpg
\end{mdcode}




\subsection{17.2 Optional media}




\begin{mdcode}
Bloque de codigo: text
13_blender_cleanup.jpg
14_unity_editor_scene.jpg
15_ui_toolkit_panel.jpg
16_shader_graph.jpg
17_github_readme_preview.jpg
18_mobile_view.jpg
19_profiler_screenshot.jpg
20_think_aloud_quote_card.jpg
\end{mdcode}




\subsection{17.3 Video assets}




\begin{mdcode}
Bloque de codigo: text
twinsight_x500_teaser_30s.mp4
twinsight_x500_demo_90s.mp4
twinsight_x500_technical_breakdown_5min.mp4
\end{mdcode}




\medskip\hrule\medskip




\section{18. Portfolio page copy blocks}




\subsection{18.1 Project card}




\begin{mdcode}
Bloque de codigo: text
TwinSight X500
Unity WebGL technical visualization for drone assembly inspection.

Interactive browser-based 3D viewer with component selection, exploded view, clipping tools, visual modes, CAD-to-realtime optimization and SUS/NASA-TLX evaluation.
\end{mdcode}




\subsection{18.2 Short homepage feature}




\begin{mdcode}
Bloque de codigo: text
TwinSight X500 converts CAD-derived drone geometry into an optimized Unity WebGL viewer for browser-based technical inspection.
\end{mdcode}




\subsection{18.3 Detailed intro}




\begin{mdcode}
Bloque de codigo: text
TwinSight X500 is my flagship real-time 3D project: a Unity WebGL technical visualization prototype designed to help users understand the assembly of a Holybro X500 V2 drone. The project combines CAD-to-realtime optimization, interactive runtime systems, technical UI, multiple visual modes and academic usability/workload evaluation.
\end{mdcode}




\subsection{18.4 Technical summary}




\begin{mdcode}
Bloque de codigo: text
Built with Unity, C#, WebGL, URP, UI Toolkit and Blender. The prototype includes component-level selection, exploded view, clipping/cross-section tools, technical information panels, multiple visual modes and optimized CAD-derived geometry for browser-based interaction.
\end{mdcode}




\subsection{18.5 Validation summary}




\begin{mdcode}
Bloque de codigo: text
The prototype was evaluated with 12 participants using SUS, NASA-TLX Raw and Think-Aloud protocols. The formative results suggested high perceived usability and lower perceived workload for the 3D viewer condition compared with 2D support in the tested tasks.
\end{mdcode}




\medskip\hrule\medskip




\section{19. GitHub README integration}




\subsection{19.1 README opening}




\begin{mdcode}
Bloque de codigo: markdown
# TwinSight X500

Unity WebGL technical visualization prototype for drone assembly inspection.

TwinSight X500 transforms CAD-derived drone geometry into an optimized browser-based 3D viewer with component selection, exploded view, cross-section tools, visual modes and technical information panels.
\end{mdcode}




\subsection{19.2 README demo block}




\begin{mdcode}
Bloque de codigo: markdown
## Demo

[![TwinSight X500 Demo](./media/twinsight-demo-thumbnail.jpg)](https://youtube.com/...)

- Live demo: [link]
- Portfolio case study: [link]
- Technical breakdown: [link]
\end{mdcode}




\subsection{19.3 README highlights}




\begin{mdcode}
Bloque de codigo: markdown
## Highlights

- Unity WebGL deployment
- CAD-to-realtime optimization
- Component-level selection
- Exploded assembly view
- Cross-section / clipping inspection
- Multiple visual modes
- UI Toolkit technical panels
- SUS + NASA-TLX evaluation
\end{mdcode}




\medskip\hrule\medskip




\section{20. ArtStation integration}




\subsection{20.1 ArtStation title}




\begin{mdcode}
Bloque de codigo: text
TwinSight X500 — Unity WebGL Technical Visualization
\end{mdcode}




\subsection{20.2 ArtStation description}




\begin{mdcode}
Bloque de codigo: text
A browser-based technical visualization prototype for inspecting the assembly of a Holybro X500 V2 drone. Built in Unity WebGL with C#, URP, UI Toolkit and Blender, the project focuses on CAD-to-realtime optimization, component selection, exploded view, cross-section tools and visual inspection modes.
\end{mdcode}




\subsection{20.3 ArtStation post order}




\begin{mdcode}
Bloque de codigo: text
01. Cover image
02. 30–60 s video
03. Short description
04. Final screenshots
05. Feature breakdown
06. CAD optimization
07. Visual modes grid
08. UI and interaction
09. Evaluation metrics
10. Tools
11. Links
\end{mdcode}




\subsection{20.4 Tags}




\begin{mdcode}
Bloque de codigo: text
Unity
Unity WebGL
Technical Art
Real-Time 3D
Technical Visualization
CAD Optimization
Drone
Interactive 3D
CSharp
URP
Blender
\end{mdcode}




Use \texttt{Digital Twin} only if the description says:




\begin{mdcode}
Bloque de codigo: text
digital-twin-adjacent visualization prototype
\end{mdcode}




not full digital twin.




\medskip\hrule\medskip




\section{21. LinkedIn integration}




\subsection{21.1 Featured item title}




\begin{mdcode}
Bloque de codigo: text
TwinSight X500 — Unity WebGL Technical Visualization
\end{mdcode}




\subsection{21.2 Featured item description}




\begin{mdcode}
Bloque de codigo: text
Interactive browser-based 3D viewer for drone assembly inspection, built with Unity, C#, WebGL, URP, UI Toolkit and Blender. Includes CAD-to-realtime optimization, component selection, exploded view, clipping tools, visual modes and SUS/NASA-TLX evaluation.
\end{mdcode}




\subsection{21.3 Launch post}




\begin{mdcode}
Bloque de codigo: text
I built TwinSight X500, a Unity WebGL technical visualization prototype for drone assembly inspection.

The project explores how interactive 3D can support spatial understanding when static 2D documentation is not enough.

Core features:
- component selection
- exploded assembly view
- cross-section inspection
- technical part panels
- multiple visual modes
- CAD-to-realtime optimization
- SUS / NASA-TLX evaluation

Stack:
Unity, C#, WebGL, URP, UI Toolkit, Blender.
\end{mdcode}




\medskip\hrule\medskip




\section{22. Interview defense notes}




\subsection{22.1 Explain in 30 seconds}




\begin{mdcode}
Bloque de codigo: text
TwinSight X500 is a Unity WebGL technical visualization prototype for inspecting a Holybro X500 V2 drone assembly. I built it to address the limitations of static 2D documentation by allowing users to inspect components, activate exploded view, cut through geometry, switch visual modes and access technical information directly in the browser. The project also includes a CAD-to-realtime optimization pipeline and formative evaluation using SUS, NASA-TLX and Think-Aloud.
\end{mdcode}




\subsection{22.2 Explain in 2 minutes}




\begin{mdcode}
Bloque de codigo: text
The project started from the problem that 2D assembly documentation requires users to mentally reconstruct spatial relationships between parts. I used a CAD-derived drone model as the basis and transformed it into a Unity WebGL prototype.

The technical work included optimizing geometry for browser use, organizing selectable components, implementing runtime interaction systems in C#, building a technical UI with UI Toolkit, adding exploded view and clipping tools, and creating several visual modes for inspection.

The project was evaluated with 12 participants using SUS, NASA-TLX Raw and Think-Aloud. The formative results suggested high perceived usability and lower perceived workload for the 3D viewer condition compared with 2D support in the tested tasks.
\end{mdcode}




\subsection{22.3 If asked about AI assistance}




\begin{mdcode}
Bloque de codigo: text
I used AI as an implementation support and iteration tool, mainly to accelerate coding, debugging and structuring parts of the work. I was responsible for the architecture, integration, testing, Unity implementation decisions, asset preparation, evaluation design and final project ownership.
\end{mdcode}




\subsection{22.4 If asked whether it is a digital twin}




\begin{mdcode}
Bloque de codigo: text
I would describe it as digital-twin-adjacent, not a full production digital twin. It represents a real physical drone assembly in interactive 3D, but it does not currently include live sensor data, operational state synchronization or predictive maintenance.
\end{mdcode}




\medskip\hrule\medskip




\section{23. Quality checklist before publishing}




\subsection{23.1 Content checklist}




\begin{mdcode}
Bloque de codigo: text
[ ] title is role-aligned
[ ] first paragraph is understandable to recruiters
[ ] technical stack visible
[ ] demo video embedded
[ ] problem clearly stated
[ ] solution clearly stated
[ ] features shown visually
[ ] CAD optimization explained
[ ] metrics included
[ ] role clearly stated
[ ] AI assistance disclosed if relevant
[ ] limitations included
[ ] links working
\end{mdcode}




\subsection{23.2 Visual checklist}




\begin{mdcode}
Bloque de codigo: text
[ ] hero image strong
[ ] UI readable
[ ] visual modes grid clear
[ ] screenshots consistent
[ ] no debug artifacts
[ ] no unfinished materials visible
[ ] captions concise
[ ] diagrams readable
[ ] video thumbnail clear
\end{mdcode}




\subsection{23.3 Claim-safety checklist}




\begin{mdcode}
Bloque de codigo: text
[ ] does not claim production deployment
[ ] does not claim full digital twin
[ ] does not claim live IoT
[ ] does not claim real thermal simulation
[ ] does not overstate evaluation causality
[ ] does not claim senior industry experience
[ ] does not hide prototype status
\end{mdcode}




\medskip\hrule\medskip




\section{24. Final recommended page skeleton}




\begin{mdcode}
Bloque de codigo: markdown
# TwinSight X500
## Unity WebGL Technical Visualization for Drone Assembly Inspection

[Hero image]
[Demo video]

TwinSight X500 is an interactive browser-based 3D viewer...

## Problem
2D documentation makes spatial relationships harder to understand...

## Solution
TwinSight converts CAD-derived drone geometry into...

## Key Features
- Component selection
- Exploded view
- Cross-section tools
- Visual modes
- Technical UI
- CAD-to-realtime optimization
- Evaluation

## Technical Pipeline
CAD → Blender → Unity → WebGL → Evaluation

## Optimization
6.5M+ source CAD routes → 95,617 optimized triangles

## Unity WebGL Architecture
Runtime interaction systems, UI Toolkit panels, visual states...

## Visual Modes
Realistic, X-ray, ghosted, blueprint, wireframe, thermal-style...

## Evaluation
SUS 91.88
NASA-TLX Raw: 8.69 viewer vs 19.89 2D support
12 participants

## My Role
Real-time 3D development, Unity WebGL integration...

## AI-Assisted Workflow
AI assistance was used as implementation support...

## Limitations
Academic/portfolio prototype. Not a full production digital twin...

## Links
- Live demo
- GitHub
- Demo video
- ArtStation breakdown
\end{mdcode}




\medskip\hrule\medskip




\section{25. Next module}




Proceed with:




\begin{mdcode}
Bloque de codigo: text
20B_portfolio_page_implementation_copy.md
\end{mdcode}




\subsection{Tool recommendation}




\begin{center}
\scriptsize
\begin{longtable}{>{\raggedright\arraybackslash}p{0.455\linewidth} >{\raggedright\arraybackslash}p{0.455\linewidth}}
\toprule
Tool & Use \\
\midrule
\endfirsthead
\toprule
Tool & Use \\
\midrule
\endhead
Agent mode & No \\
Deep Research & No \\
\bottomrule
\end{longtable}
\end{center}




Reason:




The structure is now fixed. The next file should translate it into final portfolio implementation copy: homepage, project card, project page, CTA, metadata, SEO and section-by-section text.



